\documentclass{article}

\usepackage[main, final]{neurips_2025}

\usepackage[utf8]{inputenc}
\usepackage[T1]{fontenc}
\usepackage{hyperref}
\usepackage{url}
\usepackage{booktabs}
\usepackage{amsfonts}
\usepackage{nicefrac}
\usepackage{microtype}
\usepackage{xcolor}
\PassOptionsToPackage{numbers, compress}{natbib}

\title{JEPA-NAV: A Vision-Based Joint Embedding Architecture for Safe Indoor Navigation}

\author{
  Edward Ajayi\thanks{Carnegie Mellon University Africa, Kigali, Rwanda. \texttt{eaajayi@andrew.cmu.edu}} \\
  CMU Africa, Kigali, Rwanda \\
  \texttt{eaajayi@andrew.cmu.edu}
  \And
  Emrys Lasidlaus\thanks{Carnegie Mellon University Africa, Kigali, Rwanda. \texttt{eladisla@andrew.cmu.edu}} \\
  CMU Africa, Kigali, Rwanda \\
  \texttt{eladisla@andrew.cmu.edu}
}

\begin{document}

\maketitle

\begin{abstract}
Navigating indoor environments is challenging for visually impaired individuals due to
cluttered layouts, dynamic obstacles, and partial observability. Existing vision-based
assistive systems often rely on specialized hardware, fail to generalize to unseen
environments, and rarely provide safety-aware predictions with quantified uncertainty.
We propose \textbf{JEPA-NAV}, a vision-first system that leverages Joint Embedding
Predictive Architectures (JEPA)---specifically V-JEPA~2---to learn predictive latent
representations of indoor scenes. These embeddings support affordance prediction
(walkable areas, obstacles, and risky zones) with associated uncertainty, integrated
into a risk-aware navigation planner. Our midterm establishes the V-JEPA~2 encoder
pipeline validated on ReplicaCAD, implements the full baseline (V-JEPA~2 + MLP policy
for PointNav in Habitat-Sim), and reports baseline metrics from a simple-scene run:
success rate, SPL, and collisions per episode.
\end{abstract}

% =============================================================================
\section{Introduction}
% =============================================================================

\subsection{Motivation}

Indoor navigation remains a significant challenge for visually impaired individuals
seeking independent mobility. Effective assistive systems must perceive obstacles,
identify walkable paths, and anticipate hazards while accounting for uncertainty.
Navigation becomes particularly difficult when models encounter unfamiliar
environments~\cite{najjar2022dynamic}, as traditional approaches often fail to
generalize and rarely provide safety-aware predictions.

\subsection{Problem Definition}

We address the problem of \textbf{safe, vision-based indoor navigation} that satisfies
four criteria:
\begin{enumerate}
  \item \textbf{Generalizes} to previously unseen indoor layouts without task-specific
        retraining;
  \item \textbf{Predicts} environmental affordances (walkable zones, obstacles, hazards)
        from visual input;
  \item \textbf{Quantifies uncertainty} in these predictions to support risk-aware path
        planning;
  \item \textbf{Operates} in simulation (and potentially on low-cost devices) using
        primarily RGB input, with optional depth or other modalities.
\end{enumerate}

We use \textbf{V-JEPA~2} as our visual backbone and baseline: a policy that maps
V-JEPA~2 embeddings and goal to navigation actions. Our full method, JEPA-NAV, extends
this baseline by adding affordance prediction, uncertainty modeling, and a risk-aware
planner.

% =============================================================================
\section{Related Work}
% =============================================================================

\subsection{Vision-Based Assistive Navigation}

Previous work has explored vision-based support for indoor navigation for visually
impaired users. \citet{li2018vision} proposed a system that constructs semantic indoor
maps, performs real-time obstacle detection using an RGB-D camera, and provides
multimodal feedback via audio and haptics, demonstrating improved navigation safety and
independence. \citet{na2024improving} focused on walking path generation for a robotic
SmartCane, integrating human motion constraints and motion primitives to guide users in
indoor environments. Both studies highlight the importance of real-time perception and
path planning but rely heavily on specialized hardware and system-specific
configurations.

\subsection{Predictive Visual Representations and JEPA}

Recent advances in predictive visual representations offer new opportunities for
generalizable navigation. \citet{miao_jepa-vla_2026} introduced JEPA-VLA, integrating
predictive video embeddings into vision-language-action models to capture task-relevant
temporal dynamics while ignoring unpredictable factors, improving sample efficiency and
generalization. \textbf{V-JEPA~2}~\cite{assran2025v} is a self-supervised video model
trained on large-scale video data to learn joint predictive embeddings capable of
understanding, predicting, and planning in physical environments. These representations
form the basis for our baseline and for JEPA-NAV.

\subsection{Gaps We Address}

\begin{itemize}
  \item First application of JEPA-based (V-JEPA~2) predictive embeddings for assistive
        indoor navigation in a reproducible pipeline.
  \item Uncertainty-aware affordance prediction on top of JEPA embeddings for
        risk-aware path planning.
  \item Vision-first, generalizable pipeline evaluated in standard simulators (Habitat)
        with public datasets, establishing a reproducible baseline for future research.
\end{itemize}

% =============================================================================
\section{Dataset Description}
% =============================================================================

We use \textbf{ReplicaCAD} (Habitat) for encoder validation and planned ReplicaCAD
PointNav training. ReplicaCAD is an artist-created 3D indoor dataset based on the FRL
apartment from Replica; it is designed for Habitat and supports embodied navigation and
interaction. For the midterm baseline run we also use a \textbf{simple \texttt{.glb}
scene} (Habitat test scene: \texttt{van-gogh-room.glb}) in Habitat-Sim when ReplicaCAD
requires Bullet and the full install is pending.

% \begin{table}[htbp]
% \centering
% \caption{Datasets and simulators used.}
% \label{tab:dataset}
% \begin{tabular}{lll}
% \toprule
% \textbf{Dataset / Simulator} & \textbf{Description} & \textbf{Use in Project} \\
% \midrule
% ReplicaCAD
%   & 84 scene instances, navmeshes, stages; CC BY~4.0.
%   & RGB frames; V-JEPA~2 encoder validation. \\
% Habitat-Sim
%   & 3D simulator; loads ReplicaCAD or \texttt{.glb} scenes.
%   & RGB observations; PointNav training. \\
% Simple \texttt{.glb} (\texttt{van-gogh-room})
%   & Single mesh, Habitat test scene.
%   & Midterm baseline run (no Bullet). \\
% \bottomrule
% \end{tabular}
% \end{table}

\begin{table}[htbp]
\centering
\caption{Datasets and simulators used.}
\label{tab:dataset}
\begin{tabular}{p{2.8cm} p{6cm} p{4cm}}
\toprule
\textbf{Dataset / Simulator} & \textbf{Description} & \textbf{Use in Project} \\
\midrule
ReplicaCAD
    & 84 scene instances, navmeshes, stages; CC BY~4.0.
    & RGB frames; V-JEPA~2 encoder validation. \\[4pt]
Habitat-Sim
    & 3D simulator; loads ReplicaCAD or \texttt{.glb} scenes.
    & RGB observations; PointNav training. \\[4pt]
Simple \texttt{.glb} scene
    & Single mesh, Habitat test scene.
    & Midterm baseline run (no Bullet). \\
\bottomrule
\end{tabular}
\end{table}

\noindent\textbf{Dataset statistics (midterm).}
Scenes: 84 ReplicaCAD instances (e.g.\ \texttt{apt\_0}); 80 RGB frames
($256\times256$) collected from \texttt{apt\_0} for encoder validation. Source:
Hugging~Face \texttt{ai-habitat/ReplicaCAD\_dataset}. Baseline PointNav runs reported
here use \texttt{van-gogh-room.glb} (Habitat test scene).

% =============================================================================
\section{Baseline Implementation and Results}
% =============================================================================

\subsection{Baseline Definition}

Our \textbf{baseline} is a V-JEPA~2-based navigation agent:
\begin{itemize}
  \item \textbf{Input:} RGB observation from the Habitat agent.
  \item \textbf{Encoder:} V-JEPA~2 (pretrained), producing a 1024-d latent embedding
        per frame (16-frame clip, mean-pooled).
  \item \textbf{Policy:} MLP that takes the embedding and goal (distance to target,
        angle~$\theta$) and outputs four actions: \textit{forward}, \textit{turn left},
        \textit{turn right}, \textit{stop}.
  \item \textbf{Training:} Policy gradient (actor-critic) for the PointNav task.
\end{itemize}

\noindent
There is \textbf{no} affordance head, \textbf{no} explicit uncertainty module, and
\textbf{no} risk-aware planner in the baseline. This baseline uses the same
representation that will be extended in JEPA-NAV.

% \subsection{Implementation}

% We implemented: (A)~V-JEPA~2 sanity check and encoder validation on ReplicaCAD frames;
% (B)~ReplicaCAD frame collection and V-JEPA~2 encoder statistics; (C)~full baseline via
% \texttt{train\_pointnav\_vjepa2.py} (ReplicaCAD + Bullet) and
% \texttt{train\_pointnav\_vjepa2\_simple\_scene.py} (simple \texttt{.glb} scene, no
% Bullet).

% \smallskip\noindent
% Model: \texttt{facebook/vjepa2-vitl-fpc64-256}; 16-frame clip $256\times256$
% $\to$ mean-pooled 1024-d. Policy: MLP($1024{+}2 \to 256 \to 4$). Success threshold:
% $0.36\,$m. Metrics: Success Rate (\%), SPL, Collisions per episode (mean). All results
% are logged to \texttt{logs/pointnav\_simple\_results.txt}.

\subsection{Implementation}

We implemented three components:
\begin{itemize}
    \item[(A)] V-JEPA~2 sanity check and encoder validation on ReplicaCAD frames.
    \item[(B)] ReplicaCAD frame collection and V-JEPA~2 encoder statistics.
    \item[(C)] Full baseline for both the ReplicaCAD + Bullet setting and the simple
               \texttt{.glb} scene setting (no Bullet).
\end{itemize}

\noindent
The encoder uses V-JEPA~2 with a 16-frame clip at $256\times256$ resolution,
mean-pooled to a 1024-d embedding. The policy is an MLP with architecture
$1024{+}2 \to 256 \to 4$. The success threshold is $0.36\,$m. We report
Success Rate~(\%), SPL, and mean Collisions per episode.

\subsection{Metrics}

We report three standard Habitat PointNav metrics:
\begin{itemize}
  \item \textbf{Success Rate} --- fraction of episodes where the agent reaches the goal
        within $0.36\,$m.
  \item \textbf{SPL} --- success weighted by path length.
  \item \textbf{Collisions per episode} --- mean number of collisions per episode.
\end{itemize}

\subsection{Encoder Validation Results}

V-JEPA~2 encoder output statistics:
\begin{itemize}
  \item \textbf{Phase~A (sanity):} 4 images $\to$ 16-frame clip ($256\times256$);
        output shape $(1, 2048, 1024)$; mean $0.013$, std $3.23$.
  \item \textbf{ReplicaCAD:} 80 Habitat frames $\to$ 16-frame clip; output shape
        $(1, 2048, 1024)$; mean $0.019$, std $3.03$.
\end{itemize}

\subsection{Baseline PointNav Results}

We ran the full pipeline on a simple \texttt{.glb} scene (\texttt{van-gogh-room.glb})
in Habitat-Sim with the same V-JEPA~2 + MLP policy and policy-gradient training. All
reported runs use \textbf{one scene} (\texttt{van-gogh-room}); each \textbf{episode} is
one navigation trial (one start--goal pair, up to 50~steps). Short validation runs
(4--5~episodes) and an extended run (80~episodes) completed successfully. Results are
shown in Table~\ref{tab:baseline}.

\begin{table}[htbp]
\centering
\caption{Baseline (V-JEPA~2 + policy) on simple scene. Scene and maximum steps per
  episode are fixed; only the number of training episodes varies.}
\label{tab:baseline}
\begin{tabular}{lcrrrr}
\toprule
\textbf{Run} & \textbf{Scene} & \textbf{Episodes}
  & \textbf{Success Rate (\%)} & \textbf{SPL} & \textbf{Collisions/ep} \\
\midrule
Short (4~ep)    & van-gogh-room &  4 & 0.0 & 0.000 & 9.25 \\
Tiny (5~ep)     & van-gogh-room &  5 & 0.0 & 0.000 & 9.60 \\
Extended (80~ep)& van-gogh-room & 80 & 6.2 & 0.000 & 7.88 \\
\bottomrule
\end{tabular}
\end{table}

\paragraph{Ablation: effect of number of training episodes.}
We varied the number of training episodes (4, 5, 80) while keeping the same scene
(\texttt{van-gogh-room.glb}), the same maximum steps per episode~(50), and the same
architecture. With more episodes, success rate increases ($0\% \to 6.2\%$ for
80~episodes) and mean collisions per episode decrease ($9.60 \to 7.88$), confirming
that the policy benefits from additional training. This ablation supports reporting the
80-episode run as the stronger baseline configuration.

The pipeline runs end-to-end; all metrics are computed as in Habitat PointNav.
ReplicaCAD-based training (with Habitat-Sim built with Bullet) will be run when the
cluster install completes; the same metrics and scripts apply.

% =============================================================================
\section{Methodology}
% =============================================================================

\subsection{How JEPA-NAV Extends the Baseline}

JEPA-NAV retains the same V-JEPA~2 encoder and adds three components:
\begin{enumerate}
  \item \textbf{Affordance prediction:} a head on top of V-JEPA~2 embeddings that
        predicts walkable regions, obstacles, and hazard zones (spatial map or
        per-patch labels).
  \item \textbf{Uncertainty modeling:} probabilistic or ensemble output for
        confidence/uncertainty per prediction.
  \item \textbf{Risk-aware path planning:} a planner that uses the affordance map and
        uncertainty to produce paths minimizing collision risk and preferring
        high-confidence walkable regions.
\end{enumerate}

\noindent
The agent then follows this plan or a policy conditioned on it. This design addresses
the identified gaps: generalizable predictive representations combined with explicit
affordance and uncertainty estimation for safer, assistive indoor navigation.

\subsection{Pipeline Overview}

\begin{center}
RGB (video)
$\;\to\;$ V-JEPA~2
$\;\to\;$ embeddings
$\;\to\;$ Affordance head
$\;\to\;$ walkable/obstacle/hazard map $+$ uncertainty
$\;\to\;$ Risk-aware planner
$\;\to\;$ path/cost map
$\;\to\;$ Policy/controller
$\;\to\;$ actions
\end{center}

\subsection{Planned Evaluation (Post-Midterm)}

Evaluation will use the same metrics as the baseline (Success Rate, SPL,
Collisions/ep), augmented with affordance prediction accuracy and uncertainty
calibration. We will compare JEPA-NAV against the V-JEPA~2 baseline on the same
evaluation episodes and scenes.

% =============================================================================
\section{Division of Work}
% =============================================================================

\begin{table}[htbp]
\centering
\caption{Division of work.}
\label{tab:division}
\begin{tabular}{ll}
\toprule
\textbf{Team Member} & \textbf{Contributions} \\
\midrule
Edward Ajayi
  & Literature review, V-JEPA~2 integration, training scripts, \\
  & simple-scene pipeline and experiments, report writing. \\
Emrys Lasidlaus
  & Habitat/ReplicaCAD setup, dataset documentation, baseline design and metrics, \\
  & report writing. \\
\bottomrule
\end{tabular}
\end{table}

\noindent
We split the encoder pipeline and ReplicaCAD/dataset setup between the two members,
merged efforts for the ReplicaCAD run and simple-scene baseline experiments, and
co-wrote the report.

% =============================================================================
\section{Reproducibility and Code}
% =============================================================================

Code, configs, and instructions to reproduce the baseline (and later JEPA-NAV) are
available in our repository.

\smallskip\noindent
\textbf{GitHub:} \url{https://github.com/edwardajayi/JEPA-NAV}\\
(Replace with your public repository link; if not yet public, state: ``Repository will
be made public upon course deadline.'')

\smallskip\noindent
The repository includes:
\begin{itemize}
  \item \texttt{README} with setup instructions (Habitat, V-JEPA~2, dependencies) and
        data download links;
  \item training and evaluation commands
\end{itemize}

% =============================================================================
\bibliographystyle{plainnat}
\bibliography{references}

\end{document}